\chapter*{\LARGE Abstract}
\label{ch: abstract}
\thispagestyle{plain}

Rising tail risks from extreme weather events elevate flood risk to the forefront of France’s climate governance debates. This paper quantifies the causal impact of extreme floods on labor markets and firm dynamics, thereby informing more effective private and public adaptation responses. Leveraging matched employer–employee data on all French establishments (2010–2019) and ministerial disaster declarations under the french CatNat insurance regime, I implement a Local Projection Difference-in-Differences design to address staggered treatment and dynamic responses. Flood shocks reduce employment by roughly 2\% and wages by 1\% on average in establishments affected by extreme flooding over the five years following a disaster, with losses concentrated in commerce, construction, and local services, and among small, undiversified firms. Industrial and agricultural establishments exhibit resilience. Commune-level results are more muted, possibly due to aggregate recovery dynamics, insurance transfers, and spatial reallocation of economic activity. Reduced-form estimates from this study discipline structural models of climate damages, allowing forward-looking damage assessments, and reveal heterogeneous vulnerabilities across sectors and firm structures.  \\

\noindent\textbf{Keywords:} Firm performance, Wages, Natural disasters \\
\noindent\textbf{JEL Codes:} D22, Q54, R11\\


